% ==============================================================================
% T0 Theory / FFGFT: Document 282 (Chapter Content, EN)
% Title: FFGFT in Hilbert space -- from the shared loop object to the mass operator
% Author: Johann Pascher
% Date: 16 June 2026
% Purpose: The more complete Hilbert-space translation of FFGFT. Three levels
%          (QM operations, process/time, mass). Core: the mass operator as a
%          Hermitian Z3-circulant on the C^3 fibre yielding the lepton masses as
%          a spectrum -- from two pure numbers (sqrt2, 2/9), with no fit.
%          The FFGFT<->HLV comparison is only the entry door, not the topic.
% References: Doc. 006, 172, 189, 206, 230, 231, 232, 258, 259, 270.
% Scripts: ffgft_t4_cp_divisibility.py, ffgft_t4_mass_operator.py,
%          ffgft_mass_operator_C3.py
% ==============================================================================

\section*{Abstract}
The existing Hilbert-space bridge of FFGFT (Doc.~230/231/232) already maps
previously described QM operations -- but only those, and without time. This
document makes the translation more complete and shows that it reaches the
\emph{masses}. Three levels: (0)~QM-equivalent -- both FFGFT and HLV are a
connection Laplacian on a fibre-valued Hilbert space (gauge freedom $=$ unitary
equivalence, entanglement $=$ holonomy around cycles). (1)~Process -- a
CP-divisibility test computably separates QM-internal (Markovian) from
history-coupled (non-Markovian); this is time as process structure.
(2)~Mass -- the actual new finding: the mass operator is a Hermitian
Z$_3$-circulant on the $\mathbb{C}^3$ fibre, its eigenvectors are the three
Z$_3$ characters (the generations), its spectrum $S^2$ are the masses. With two
pure numbers -- $r=\sqrt2$ (Koide $Q=2/3$) and the angle $\theta=2/9$ -- it
reproduces $m_\mu/m_e$ and $m_\tau/m_\mu$ to four--five digits, without the
formerly fitted ladder coefficient. Here $2/9=2/3^2$ is not a tuned decimal but
a structural number on the footing of $\xi$: a foundational condition of the
fractal open space. The FFGFT$\leftrightarrow$HLV comparison is only the entry
door; no theoretical equivalence is claimed.

% ==============================================================================
\section{Occasion and scope}
% ==============================================================================

The exchange with Marcel Kr\"uger (HLV/HICE) raised the question whether FFGFT
and HLV are connected. Pursuing it surfaced the essential point: FFGFT's
Hilbert-space mapping is \emph{incomplete} -- it translates the
already-described QM operations, but \emph{time} does not appear in it, although
it very much exists in the real world and in the $T^4$ tori. And as long as
time is missing, the \emph{masses} are missing. This document closes that gap
as far as is honestly possible at present.

\begin{important}[Scope rule]
\emph{No} theoretical equivalence FFGFT~$=$~HLV is claimed. The HLV comparison
only supplies the shared object (connection Laplacian) and the process
observable (non-Markovianity). The topic is the completion of FFGFT's own
Hilbert-space translation up to the masses.
\end{important}

% ==============================================================================
\section{The Hilbert space explicitly}
% ==============================================================================

Before computing, the space itself. FFGFT's full Hilbert space is the tensor
product of base and fibre:
\begin{equation}
  \mathcal{H}\;=\;L^2(T^4)\;\otimes\;\mathbb{C}^3 .
\end{equation}

\begin{notationBox}[Bases and inner product]
\textbf{Base $L^2(T^4)$} -- functions on the four-torus, two natural bases:
\begin{itemize}
  \item node basis $\ket{x}$, $x\in T^4$ (nodes of the discrete torus graph);
  \item winding basis $\ket{n}$, $n=(n_1,n_2,n_3,n_4)\in\mathbb{Z}^4$, since
        $H_1(T^4)=\mathbb{Z}^4$ (Doc.~189) -- eigenstates of the translation
        generators, the $n_i$ being the winding numbers.
\end{itemize}
\textbf{Fibre $\mathbb{C}^3$} -- the three Z$_3$-symmetric ICE channels
(Doc.~206), two bases:
\begin{itemize}
  \item channel basis $\ket{c}$, $c=1,2,3$;
  \item Z$_3$ character basis
        $\ket{f_k}=\tfrac{1}{\sqrt3}\sum_c\omega^{kc}\ket{c}$, $k=0,1,2$,
        $\omega=e^{2\pi i/3}$ -- the \emph{generations}.
\end{itemize}
\textbf{Inner product:} $\langle\psi|\phi\rangle=\sum_{x\in T^4}\sum_c
\overline{\psi(x,c)}\,\phi(x,c)$. A general state:
$\ket{\psi}=\sum_{n,k}a_{n,k}\,\ket{n}\otimes\ket{f_k}$.
\end{notationBox}

Three operator families act on this space -- and they split the physics
cleanly:

\begin{technical}[The three operators on $\mathcal{H}$]
\textbf{(1) Connection Laplacian $L$ on the base} (gauge/holonomy sector):
$L=D-\sum_{\langle vw\rangle}U_{vw}$, $U_{vw}=e^{iA_{vw}}$. Eigenvalues $=$
internal mode numbers; gauge transformation $=$ unitary conjugation; holonomy
around cycles $=$ entanglement (Doc.~230). Lives on $L^2(T^4)$.\\[2pt]
\textbf{(2) Time generator $H_t$} (energy/mass sector): generator of
translation in the $\tilde T$ direction of the torus. Via $\tilde T\cdot m=1$ its
eigenvalue is the rest mass; it sets the scale $M=1/\tilde T$.\\[2pt]
\textbf{(3) Mass operator $S$ on the fibre} (generation sector): the Hermitian
Z$_3$-circulant on $\mathbb{C}^3$. Eigenvalues $\sqrt{m_k}$, eigenvectors $=$
Z$_3$ characters $=$ generations.
\end{technical}

Thus the physics splits: the base $L^2(T^4)$ carries gauge structure,
entanglement and dynamics; the fibre $\mathbb{C}^3$ carries generation and mass;
the $\tilde T$ direction sets the scale. The full operator of a sector is a
product (base operator)\,$\otimes$\,(fibre Z$_3$-circulant).

% ==============================================================================
\section{Level 0 -- the shared object: connection Laplacian}
% ==============================================================================

At its core the QM mapping is a connection Laplacian on a fibre-valued function
space over a discrete base -- the $T^4$ for FFGFT, the cell graph for HLV:
\begin{equation}
  L \;=\; D \;-\!\!\sum_{\langle vw\rangle} U_{vw}\,,
  \qquad U_{vw}=e^{iA_{vw}}\in \mathrm{U}(n),
  \qquad \Omega(C)=\textstyle\sum_{e\in C}A_e \pmod{2\pi}.
\end{equation}

\begin{keyresult}[What FFGFT achieves at this level]
States $=$ fibre-valued fields; QM operation $=$ unitary parallel transport
$U_{vw}$; gauge freedom $=$ unitary equivalence (gauge-null, passed on both
sides, $\sim10^{-14}$); entanglement $=$ holonomy around cycles (Doc.~230) with
$H_1(T^4)=\mathbb{Z}^4$ (Doc.~189); CP phase $=$ accumulated torsion phase
(Doc.~172). Counterparts: FFGFT loop Hamiltonian (Doc.~231/232)
$\leftrightarrow$ HLV $L_{\mathrm{HICE}}$.
\end{keyresult}

This is solid, but it is only QM cast as a lattice gauge model -- and it
contains neither time nor mass.

% ==============================================================================
\section{Level 1 -- the process level: CP-divisibility}
% ==============================================================================

The first extension beyond static QM is time as \emph{process}. A fibre qubit
is dephased by a bath whose frequencies $\omega_k=\sqrt{\lambda_k}$ are the
loop frequencies of the $T^4$ connection Laplacian (independent boson):
\begin{equation}
  \Gamma(t)=\sum_k\frac{2g_k^2}{\omega_k^2}\bigl(1-\cos\omega_k t\bigr),\quad
  c(t)=e^{-\Gamma(t)},\quad
  \gamma(t)=\dot\Gamma(t)=\sum_k\frac{2g_k^2}{\omega_k}\sin\omega_k t.
\end{equation}
Two reduced descriptions of the same operator are compared: (M) memoryless /
time-local (constant rate $\Rightarrow$ CP-divisible) and (NM) history-coupled
(full $T^4$ bath $\Rightarrow$ may violate CP-divisibility). Two standard
witnesses (BLP trace-distance back-flow, RHP/Choi):

\begin{center}
\begin{tabular}{lcccc}
\toprule
Description & BLP back-flow & $D$ revives & min Choi eig & $\gamma<0$ \\
\midrule
(M)\ \ memoryless / time-local      & $0.000$ & no  & $+0.013$ & --- \\
(NM) history-coupled ($T^4$)        & $5.125$ & yes & $-0.041$ & 1928 \\
\bottomrule
\end{tabular}
\end{center}

Only the history-coupled variant leaves the Markovian class. This is the
operational \enquote{beyond-Markov} signature -- but beyond the time-local
\emph{reduced} description, not beyond quantum theory.

% ==============================================================================
\section{Level 2 -- time and mass: the Hilbert-space mass operator}
% ==============================================================================

Here lies the actual finding. A mass is a rest energy: the eigenvalue of the
\emph{time-translation generator} (Casimir $m^2=E^2-p^2$). The basic relation
$\tilde T\cdot m=1$ makes mass the \emph{dual} of the time cycle. A static graph
Laplacian has discarded time -- its eigenvalues are internal mode numbers, not
masses. Bringing time back yields masses; and since $\hbar,c$ are merely SI
conversion factors, the physical content is the \emph{ratio}, while the scale
$M=1/\tilde T$ is set by the time cycle.

\begin{warningboxenv}[A plain Laplacian is not enough]
A symmetric $T^4$ connection Laplacian with $m^2=\lambda_k$ gives the three
lowest \enquote{masses} in the ratio $1{:}1{:}1.46$ and Koide $Q=0.345$ --
O(1) ratios, no hierarchy, Koide far from $2/3$. Time as a uniform extra
dimension does not suffice; the Z$_3$ channel structure is required.
\end{warningboxenv}

\subsection{The operator}

The fibre is $\mathbb{C}^3$ -- the three Z$_3$-symmetric ICE channels (Doc.~206:
$\det A_\xi=1-\xi$), i.e.\ the generation space. The full space is
$L^2(T^4)\otimes\mathbb{C}^3$; the $\tilde T$ cycle of the torus sets the scale.
The mass operator is a \emph{Hermitian Z$_3$-circulant}
\begin{equation}
  S=\mathrm{circ}(t_0,t_1,t_2),\qquad t_0=M\in\mathbb{R},\quad t_2=\bar t_1,
\end{equation}
with eigenvalues
\begin{equation}
  \mu_k = t_0 + 2|t_1|\cos\!\bigl(\arg t_1 + \tfrac{2\pi k}{3}\bigr)
        = M\bigl(1 + r\cos(\theta + \tfrac{2\pi k}{3})\bigr),
  \qquad k=0,1,2,
\end{equation}
where $r=2|t_1|/M$ and $\theta=\arg t_1$. The \emph{eigenvectors} are the three
Z$_3$ Fourier characters $f_k=(1,\omega^k,\omega^{2k})/\sqrt3$,
$\omega=e^{2\pi i/3}$ -- numerically verified with overlap $1.0000$. The three
generations \emph{are} the Z$_3$ characters; the masses are the spectrum of
$S^2$.

\subsection{Koide from the operator}

With $\sqrt{m_k}=\mu_k$ the Koide identity follows directly:
\begin{equation}
  Q=\frac{\sum_k m_k}{\bigl(\sum_k\sqrt{m_k}\bigr)^2}=\frac{1+r^2/2}{3}
  \;\Longrightarrow\; Q=\tfrac23 \iff r=\sqrt2.
\end{equation}
Koide $Q=2/3$ is thus nothing but the operator condition $r=\sqrt2$ -- the
Z$_3$+$\sqrt2$ structure (Doc.~258/259), not a fit.

\subsection{The lepton masses from two pure numbers}

With $r=\sqrt2$ and $\theta=2/9$ -- two pure numbers, no fitted coefficient:

\begin{center}
\begin{tabular}{lccc}
\toprule
 & operator & data & error \\
\midrule
$m_\mu/m_e$    & $206.770$ & $206.768$ & $0.00\,\%$ \\
$m_\tau/m_\mu$ & $16.818$  & $16.817$  & $0.01\,\%$ \\
Koide $Q$      & $0.666667$ & $0.666661$ & exactly $2/3$ \\
\bottomrule
\end{tabular}
\end{center}

\begin{keyresult}[More than a fit]
The formerly fitted ladder coefficient ($r_e\approx1.349$, Doc.~006) has
\emph{vanished} -- replaced by the structural number $2/9$. And it is more than
an adjustment: \emph{one} angle $\theta$ reproduces \emph{both} independent
ratios simultaneously (other $\theta$ break $m_\mu/m_e$ at once). Two pure
numbers $\to$ three masses to 4--5 digits.
\end{keyresult}

\begin{warningboxenv}[Numerical caution: the electron is a cancellation]
The lightest eigenvalue is a near-total cancellation,
$\mu_e = M\bigl(1+r\cos(\theta+\tfrac{2\pi}{3})\bigr)
       = M(1-0.95965) = 0.04035\,M$ -- the difference of two nearly equal
numbers. Rounding the cosines to only three digits shifts $\mu_e$ by orders of
magnitude: taking $\cos(\theta+\tfrac{2\pi}{3})$ as $-0.658$ instead of the
correct $-0.6786$ (a $3\,\%$ error) turns $\mu_e=0.040$ into $0.069$ (a
$72\,\%$ error), and since $m_e=\mu_e^2$ this flips $m_\mu/m_e$ from $207$ to
$63$. The ratios must therefore be checked at \emph{full machine precision};
precisely because the electron is the lightest state, $m_\mu/m_e$ is the
sharpest test of whether $\theta$ is exactly $2/9$ (script
\texttt{ffgft\_mass\_operator\_C3.py}, numpy precision).
\end{warningboxenv}

\emph{How precise?} Against the high-precision $e/\mu$ data the angle is
confirmed to $|\theta-2/9|\approx1.8\cdot10^{-7}$ -- seven significant figures.
The lepton masses are among the most precisely measured quantities in physics;
the data will not get sharper. The empirical question is therefore settled:
$2/9$ is the number.

\subsection{$2/9$ as a foundational condition}

Must $2/9$ be derived from something deeper? No. The decisive distinction is
between a \emph{fitted continuous knob} (bad -- that was $r_e$) and a
\emph{fixed structural rational} forced by the geometry (fine -- like $2/3$,
like $2\pi/3$, like the three-channel count itself).

\begin{foundation}[Structural number, not a fit]
$2/9=2/3^2$ is a clean rational, not a tuned decimal -- exactly the kind of
primitive the fractal open space can carry, \emph{on the footing of $\xi$}.
Like $\xi=4/30000$ it needs no deeper derivation; it is a foundational
condition of the geometry.
\end{foundation}

This shifts the leptons from a \emph{$\xi$-power} story to a \emph{Z$_3$
geometry} story: the ratios hinge on the pure numbers $\sqrt2$ and $2/9$, not on
$\xi$ (which does not appear in the ratio). This fits the fact that the
$\xi$-ladder was the weak part for leptons anyway. The fractal space then
carries a small set of primitives: $\xi$ for the coupling sector,
$(\sqrt2,\,2/9)$ for the generation sector.

% ==============================================================================
\section{Computing in the Hilbert space: recipe and examples}
% ==============================================================================

The procedure is always the same -- every physical quantity is an eigenvalue,
an eigenvalue relation, or an overlap of eigenbases on $\mathcal{H}$.

\begin{keyresult}[The recipe -- four steps]
\begin{enumerate}
  \item \textbf{Choose the sector} -- which quantity (masses, a process
        observable, a mixing angle, \dots).
  \item \textbf{Write the operator} -- the corresponding self-adjoint operator
        on $\mathcal{H}$ (Z$_3$-circulant on $\mathbb{C}^3$ for masses;
        time-extended connection Laplacian for dynamics; overlap of two
        eigenbases for mixing).
  \item \textbf{Diagonalise} -- eigenvalues $=$ quantities (as dimensionless
        ratios), eigenvectors $=$ states; relations (Koide) follow
        algebraically.
  \item \textbf{Scale $+$ SI} -- scale $M=1/\tilde T$ from the time cycle, then
        $\hbar,c$ as pure SI conversion.
\end{enumerate}
\end{keyresult}

\paragraph{Already computed.}
\begin{itemize}
  \item \emph{Masses:} $S$ on $\mathbb{C}^3$, eigenvalues$^2=$ masses, ratios
        from $(\sqrt2,2/9)$ (Level~2).
  \item \emph{Koide:} $Q=(1+r^2/2)/3$ from the eigenvalue structure, $=2/3$ at
        $r=\sqrt2$ (Level~2).
  \item \emph{Process / non-Markov:} the time-extended operator on
        $L^2(T^4)\otimes(\text{qubit})$, CP-divisibility witness (Level~1).
\end{itemize}

\paragraph{How further quantities are computed (programme).}
The same recipe, new operators -- the method is fixed, the results are still
open:
\begin{itemize}
  \item \textbf{Mixing angles (PMNS/CKM-like).} If two sectors (say charged
        leptons and neutrinos) are each Z$_3$-circulants, but in fibre bases
        \emph{rotated} against one another, then the mixing matrix is the
        overlap $\langle f_k^{(A)}|f_l^{(B)}\rangle$ of their eigenvectors. The
        mixing angles encode the \emph{rotation angle} between the two Z$_3$
        structures -- a single geometric number per sector pair, computed just
        like $2/9$.
  \item \textbf{Further sectors (quarks).} The same scheme: one Z$_3$-circulant
        per charge sector with its own structural angle; the open question is
        the same -- which rationals does the topology force.
  \item \textbf{Dynamics and correlations.} Holonomy and entanglement
        observables from the base connection Laplacian (Doc.~230), already
        demonstrated in the witness scripts.
\end{itemize}

\begin{important}[What \enquote{computing} means here]
Every quantity is an eigenvalue, an eigenvalue relation, or an eigenbasis
overlap on $\mathcal{H}$, fixed by structural rationals plus the single scale
$\tilde T$. That is the sense in which the Hilbert-space mapping
\enquote{computes further things} -- not through new fits, but through new
operators following the same recipe. The quark sectors are computed as an
operator spectrum (next section) -- faithfully, but without pure numbers; the
mixing angles remain open. What is fixed is the method.
\end{important}

% ==============================================================================
\section{The three sectors explicitly: faithful translation, lepton-specific numbers}
% ==============================================================================

Because the Hilbert-space mapping is a \emph{faithful translation} (the
$T^4\!\leftrightarrow$\,Hilbert bijection, Doc.~270, type~III), it must also
carry the masses FFGFT obtains via the $\xi$-route ($m=r\,\xi^p\,v$, Doc.~006)
-- leptons \emph{and} quarks. Here that is carried out: each charge sector is
exactly the spectrum $S^2$ of a Hermitian Z$_3$-circulant.

\begin{center}
\begin{tabular}{lcccc}
\toprule
Sector & $M$ & $r$ & Koide $Q$ & max.\ rel.\ error \\
\midrule
Leptons $(e,\mu,\tau)$ & $17.72$ & $\sqrt2=1.4142$ & $0.6667$ & $2\cdot10^{-14}$ \\
Up $(u,c,t)$           & $150.7$ & $1.757$         & $0.848$  & $7\cdot10^{-14}$ \\
Down $(d,s,b)$         & $25.73$ & $1.546$         & $0.732$  & $2\cdot10^{-14}$ \\
\bottomrule
\end{tabular}
\end{center}

In every sector $S^2$ reproduces the masses to machine precision (script
\texttt{ffgft\_z3\_alle\_sektoren.py}). The Hilbert-space calculation is thus
\emph{universal} -- the quarks are included.

\begin{warningboxenv}[A faithful translation is no gain by itself]
This universality is \emph{automatic}: the circle map
$\sqrt{m_j}=M\bigl(1+r\cos(\theta+\tfrac{2\pi j}{3})\bigr)$ is a
\emph{bijection} (three masses $\leftrightarrow$ $M,r,\theta$). Any triplet can
be written into it exactly -- it \emph{must}, otherwise it would not be a
translation. \enquote{It works in Hilbert space} is therefore not a prediction
but the meaning of faithfulness (cf.\ the $\xi^N/\varphi^N$ triviality,
Doc.~190 P35).
\end{warningboxenv}

The genuine content lies in whether the operator parameters are \emph{pure
numbers}:
\begin{itemize}
  \item \textbf{Leptons:} $r=\sqrt2$ -- which \emph{is} Koide $Q=2/3$, the
        Z$_3+\sqrt2$ structural condition (Doc.~258/259) -- \emph{and}
        $\theta=2/9$ (a clean rational). Two pure numbers (see above).
  \item \textbf{Quarks:} $r_{\text{up}}\approx1.76$,
        $r_{\text{down}}\approx1.55$ -- \emph{not} $\sqrt2$; the angles are not
        clean rationals. The operator carries the masses faithfully, but with
        \emph{un-special} parameters.
\end{itemize}

This matches the $\xi$-route point for point (Doc.~006): there the lepton ratios
carry real structure (Koide), while the quark $(r,p)$ are per-particle
classification. A faithful translation \emph{must} agree on \emph{where} the
structure is -- and it does.

\begin{keyresult}[Consequence for a universal Z$_3$ angle]
The attempt to impose the lepton condition $r=\sqrt2$ (i.e.\ $k=1$, Koide $2/3$)
on the quark sectors is \emph{falsified} by the faithful translation itself:
the up sector needs $r\approx1.76$, the down sector $r\approx1.55$. There is no
room for $\sqrt2$ -- and hence none for a single universal $k=1$ Z$_3$ angle
across sectors.
\end{keyresult}

% ==============================================================================
\section{The same scale carries the couplings too: $\alpha$ and (briefly) $G$}
% ==============================================================================

The Z$_3$ operator of the lepton sector fixes not only the masses but also a
\emph{scale} $M_T$ and the dimensionless eigenvalues $\mu_e,\mu_\mu,\mu_\tau$.
The fine-structure constant uses exactly these quantities -- it is therefore not
a separate input.

\subsection{$\alpha$ as an eigenvalue relation on $\mathcal H$}

The very recipe -- every quantity an eigenvalue or eigenvalue relation on
$\mathcal H$ -- yields $\alpha$ too, from the \emph{same} operator that gives the
masses. The eigenvalues of $S$ are $\lambda_j=\sqrt{m_j}$; the fine-structure
constant is the product of the two smaller ones, scaled by $\xi$:
\begin{equation}
  \alpha=\xi\left(\frac{\lambda_e\,\lambda_\mu}{\mathrm{MeV}}\right)^2,
  \qquad \lambda_e\,\lambda_\mu=\sqrt{m_e\,m_\mu}=E_0 .
\end{equation}
This is the $\xi$-route ($\alpha=\xi(E_0/\mathrm{MeV})^2$, Doc.~011), but as a
\emph{Hilbert-space variant}: $E_0$ is not a separate input but an eigenvalue
relation of the same operator built from $\sqrt2$ and $2/9$. Thus $\alpha$ and
the lepton masses follow from \emph{one} object and \emph{one} scale.
Numerically: the operator/geometric value gives $1/\alpha=138.9$; the corpus
tunes $E_0=7.398$~MeV to $1/\alpha=137.04$ -- the final near-one-percent sits in
the precise value of $E_0$ (radiative-correction level). Script
\texttt{ffgft\_kopplungen\_hilbertraum.py}.

\subsection{$G$ -- structurally through $\xi$, involved only in the SI conversion}

The gravitational coupling enters \emph{likewise} through $\xi$, via the T0
fundamental relation
\begin{equation}
  \xi=2\sqrt{G\,m}\;\Longrightarrow\; G=\frac{\xi^2}{4m}
\end{equation}
(Doc.~006, calc-o). The structural finding is the same as for $\alpha$: the
coupling is \emph{not independent} but fixed by $\xi$. The Hilbert-space mapping
supplies the scale here: $m$ is the operator scale $M_T=M^2=1/\tilde T$ of the
lepton sector. Unlike the dimensionless $\alpha$, however, $G$ carries
dimensions; its SI realisation requires the full conversion through
$\ell_P,c,\hbar$ and is therefore considerably more involved -- it is carried out
in calc-o and yields $G_{\text{SI}}$ at the per-mille level. What is involved is
the \emph{SI conversion}, not the structural origin.

\begin{keyresult}[Masses and couplings from one object]
$\alpha$ and the lepton masses follow from the \emph{same} Z$_3$ operator and its
scale; the gravitational coupling is fixed by the \emph{same} $\xi$. The
pure-number content sits in $\xi$ and the operator structure; the SI numerical
values -- $G$ especially -- are conversions, not additional inputs.
\end{keyresult}

% ==============================================================================
\section{The QM translation and what its extension changes}
% ==============================================================================

The translation of the QM \emph{operations} has its own document
(Doc.~230/231/232): gauge freedom $=$ unitary equivalence, entanglement $=$
holonomy around cycles, Bell correlations $=$ topological connections on $T^4$
(Doc.~230). That translation does \emph{not} contain mass and time -- it maps the
operations only. The natural question: does anything in the QM results change when
it is extended -- by the $\mathbb{C}^3$ mass fibre (Level~2) and time as process
(Level~1)? Computed (\texttt{ffgft\_qm\_uebersetzung\_erweiterung.py}):

\begin{center}
\begin{tabular}{lcc}
\toprule
Extension & QM result & change \\
\midrule
mass ($\mathbb{C}^3$ fibre) & CHSH/Bell on the base & $4\cdot10^{-16}$ (none) \\
time (process) & coherence back-flow & $0 \to 0.53$ \\
time (process) & min.\ Choi eig & $+0.0003 \to -0.0056$ \\
\bottomrule
\end{tabular}
\end{center}

\textbf{Mass} acts on a \emph{separate} tensor factor. A QM correlation (CHSH on
two base qubits, $=2\sqrt2$, Tsirelson) is identical to within $4\cdot10^{-16}$
with and without the $\mathbb{C}^3$ mass fibre attached. The mass extension only
enlarges the space and supplies the masses; \emph{no} static QM result changes.

\textbf{Time} changes the \emph{dynamical} quantities. The QM bridge treats the
reduced dynamics as time-local (unitary/Markovian); the extended translation
carries the loop history. The same fibre-qubit coherence $c(t)$ is Markovian and
CP-divisible with no back-flow, but history-coupled it is non-Markovian:
back-flow $0.53$, the min.\ Choi eigenvalue turns negative,
$\max|\Delta c(t)|\approx0.78$. Single-time quantities remain; sequential/dynamical
ones acquire the non-Markovian correction.

\begin{keyresult}[Slightly incomplete -- only in scope]
The QM translation is not wrong but \emph{incomplete in scope}: no mass,
time-local. Extended, it leaves every \emph{static} QM result unchanged (mass
factorises) and adds exactly \emph{one} new effect -- non-Markovian memory in the
dynamics -- which is itself still within quantum theory (a traced-out unitary bath
reproduces it, Level~1). The expectation that the extension \enquote{has an
influence on the results} therefore holds, but precisely located: on \emph{time},
not mass, and only on \emph{dynamical} quantities.
\end{keyresult}

\smallskip
\noindent\textbf{Test duration decides.} The difference lives entirely in time and
has two scales: a \emph{fast} decoherence time (coherence below $1/e$ at
$t\approx0.4$, nearly gone at $t\approx1.4$) and a \emph{late} memory/recurrence
time (first revival at $t\approx2.8$, about $6\times$ later). Hence three regimes:
(i)~very short test ($t\ll0.4$) -- coherence $\approx1$, QM-only and extended
translation identical, the QM result exact; (ii)~into the decoherence -- both fall
essentially together, the Markovian description suffices; (iii)~only from the
recurrence time on do the curves separate (Markovian $\sim0.53$ vs non-Markovian
$\sim0.22$). Since a single operation typically ends \emph{before} decoherence,
the non-Markovian term is \emph{operationally invisible} to fast operations; it
becomes measurable only once the test duration reaches the T$^4$ memory time -- in
deliberately long, coherence-preserving, sequential experiments.

\smallskip
\noindent\textbf{Does it scale with many qubits or with mass?} No -- least of all
there. Massive systems couple more strongly to the bath and therefore decohere
\emph{faster}: at four-fold coupling the decoherence time drops from
$t\approx0.44$ to $0.11$, and the absolute coherence at the first revival falls
from $0.22$ to $\sim10^{-5}$ -- the non-Markovian term is then \emph{buried}. Mass
thus makes it less visible, not more; what remains is faster \emph{ordinary}
(Markovian) decoherence. With many thousands of qubits it is not the memory term
that grows but ordinary decoherence and error accumulation (gate errors,
crosstalk) -- Markovian and addressed by error correction; whether revivals
survive at all depends on the bath \emph{structure} (incommensurate frequencies
push the Poincar\'e recurrence into the practically unreachable), not on the qubit
count. The one legitimate caveat is \emph{correlated} (non-Markovian) noise across
gates -- relevant only if the bath memory time approaches the sequence duration,
which fast gates and fast-decohering (heavy) systems precisely prevent.

% ==============================================================================
\section{Compact or recursive? When the compact computation suffices}
% ==============================================================================

One question decides how far the compact Hilbert-space computation reaches:
does the compact (cyclic, finite, one-shot) variant \emph{always} suffice --
so that decompactification is only the final step to SI -- or are there
quantities computable \emph{only recursively}? The decisive criterion is
\textbf{linearity}.

\begin{keyresult}[Linear $\Rightarrow$ compact suffices]
If the dynamics is generated by a \emph{fixed linear} operator on $\mathcal{H}$
with a discrete (compact) spectrum, the closed spectral representation always
exists:
\begin{equation}
  \langle A(t)\rangle=\sum_{k,l}c_{kl}\,e^{i(\omega_k-\omega_l)t}.
\end{equation}
One diagonalisation suffices. The memory kernel $K(t)$, which looks
\emph{recursive} in the time representation (convolution, memory), is merely the
Fourier transform of the discrete spectral density
$\sum_k g_k^2\,\delta(\omega-\omega_k)$. The recursion is thus only \emph{one
representation} (time domain); the spectral sum is the equivalent closed one.
For a discrete spectrum the closed form always exists -- decompactification
happens only at the end, at the SI step.
\end{keyresult}

Concrete evidence: the CP-divisibility test (Level~1) obtained the
\emph{non-Markovian} result -- back-flow, revivals -- from exactly such a finite
spectral sum, in closed form, \emph{without} recursion. Even \enquote{memory}
was compact-computable.

\begin{important}[When recursion is forced]
Recursion becomes unavoidable only if one of the two linearity prerequisites
breaks:
\begin{enumerate}
  \item \textbf{Fundamentally continuous spectrum} -- then no finite sum. But
        the continuum appears only at the SI/experienced level; fundamentally
        it is compact. This case drops out.
  \item \textbf{Nonlinear / self-referential dynamics} -- the generator depends
        on the \emph{state or the history itself}, not only on time. Then there
        is no fixed eigenbasis, no diagonalisation, and the recursion is
        genuine.
\end{enumerate}
\end{important}

The question thereby reduces to a \emph{classification}: for each quantity ask
-- is the generator a fixed linear operator on $\mathcal{H}$? Yes $\Rightarrow$
compact, one-shot, SI only at the end. Does it depend on the state/trajectory
$\Rightarrow$ recursion possibly irreducible.

Mass sector, gauge, entanglement and mixing are linear -- for them the compact
computation suffices, and HLV's recursions are \emph{avoidable} there. The only
candidate for genuine recursion is the \emph{self-referential} level (Doc.~248,
§8: the origin thought as a material $T^4$ process): where the process
processes itself, the generator could become state-dependent. To check this --
linear or self-referential, quantity by quantity -- is a well-defined programme
and marks the single place where the open/recursive variant would do real,
non-eliminable work.

% ==============================================================================
\section{Claim-state and limits}
% ==============================================================================

\begin{important}[What is shown -- and what is not]
\begin{enumerate}
  \item \textbf{A genuine Hilbert-space mapping.} The mass operator is a
        concrete Hermitian operator on $\mathbb{C}^3$ with eigenvectors
        (generations) and spectrum (masses) -- not a parametrisation.
  \item \textbf{Koide forward.} $Q=2/3$ follows exactly from $r=\sqrt2$;
        structure, not a fit.
  \item \textbf{Two numbers, three masses.} $(\sqrt2,2/9)$ reproduce both
        ratios to 4--5 digits; one angle covers two observables (a genuine
        over-determination).
  \item \textbf{Structural assumptions.} The \emph{form} (Z$_3$-circulant) and
        $r=\sqrt2$ are claimed geometric (Doc.~206/258/259).
  \item \textbf{Empirically exhausted.} $\theta=2/9$ is confirmed against the
        $e/\mu$ data to $\sim2\cdot10^{-7}$ (seven figures in the angle); the
        ratios agree to $10^{-5}$--$6\cdot10^{-5}$, for the $\tau$ within the
        measurement uncertainty. More precise lepton data are not to be
        expected -- the empirical question is answered. The tiny residual
        ($Q=0.6666605$ vs.\ $2/3$) sits at the level of radiative corrections,
        not a defect of the identification.
  \item \textbf{The only open question is theoretical:} is $\theta=2/9$ (like
        $r=\sqrt2$) forced by the $T^4$/Z$_3$ topology?
  \item \textbf{Not beyond quantum theory.} The CP test (Level~1) shows
        non-Markovianity in the reduced dynamics; a traced-out unitary bath
        reproduces it.
\end{enumerate}
\end{important}

The sharp, now reachable remainder is thereby precise: not \enquote{derive a
messy number}, but the single clear question -- \textbf{is $\theta=2/9$ forced
by the $T^4$/Z$_3$ topology, as $2/3$ is?} A single rational, not a ladder fit.

% ==============================================================================
\section{Conclusion}
% ==============================================================================

\begin{conclusionBox}[Conclusion]
FFGFT's Hilbert-space translation reaches further than previously shown.
Level~0 maps the QM operations cleanly (shared with HLV). Level~1 brings time
as process (CP-divisibility, non-Markovianity). Level~2 brings time as geometry
and hence the \emph{masses}: the mass operator is a Hermitian Z$_3$-circulant
on the $\mathbb{C}^3$ fibre whose spectrum yields the lepton masses from two
pure numbers $\sqrt2$ and $2/9$ -- Koide exact, both ratios to 4--5 digits,
without the old fitted coefficient. $2/9=2/3^2$ is a structural number on the
footing of $\xi$, a foundational condition of the fractal space. The whole
lepton problem is thereby reduced to a single geometric question: whether
$T^4$/Z$_3$ forces the angle $2/9$.
\end{conclusionBox}

\bigskip
\noindent\textbf{Closing remark: translation is not justification.}
Simplified, the result of these extensions fits in one sentence: \emph{a faithful
translation into Hilbert-space notation is possible for all masses and all
constants -- but precisely because it is faithful, it provides no
justification.} Faithfulness is mathematically a bijection: three masses
correspond invertibly and uniquely to the triple $(M,r,\theta)$, $\alpha$ is an
eigenvalue relation of the same operator, $G$ follows through the same scale
$1/\tilde T$. A bijection, however, carries \emph{everything} over and explains
\emph{nothing}; it is a change of coordinates, not a derivation. That a quantity
\enquote{can be computed in Hilbert space} is therefore never a finding in its
own right -- it is the meaning of faithfulness.

\smallskip
\noindent
The justification comes from the original $T^4$ torus representation, which is
generatively prior (Doc.~206/270): the Hilbert space is the \emph{image}, the
torus the \emph{original}. Only there are $\sqrt2$ (the Koide condition $Q=2/3$),
the angle $2/9$ and the primitive $\xi$ \emph{geometric resp.\ topological facts}
rather than read-off parameters. The clean division of labour is thus: the
Hilbert space \emph{computes and represents} (universally -- leptons, quarks,
$\alpha$, $G$); the $T^4$ \emph{justifies and derives}. Even where the
translation yields \emph{no} pure numbers -- for the quarks -- this is not a
defect of the translation but an expression of the fact that the torus geometry
forces none there. The open question of the whole construction is therefore a
single one, and it must be posed on the torus side: whether $T^4$/Z$_3$ forces
the pure numbers $\sqrt2$, $2/9$ and $\xi$.

\bigskip
\noindent\textbf{Accompanying scripts (numpy).}
\texttt{ffgft\_t4\_cp\_divisibility.py} (Level~1, CP-divisibility),
\texttt{ffgft\_t4\_mass\_operator.py} (Z$_3$-circle analysis, KK comparison),
\texttt{ffgft\_mass\_operator\_C3.py} (the mass operator on $\mathbb{C}^3$),
\texttt{ffgft\_z3\_alle\_sektoren.py} (all three sectors as an operator spectrum),
\texttt{ffgft\_kopplungen\_hilbertraum.py} ($\alpha$ and $G$ via the Hilbert-space
mapping),
\texttt{ffgft\_qm\_uebersetzung\_erweiterung.py} (mass/time extension of the QM
translation).
